\chapter*{Appendix D: Suture Materials \& Selection Guide}
\addcontentsline{toc}{chapter}{Appendix D: Suture Materials \& Selection Guide}
\markboth{Appendix D: Suture Materials}{Appendix D: Suture Materials}

The selection of appropriate suture material is a critical component of surgical technique. Sutures are chosen based on the tissue's healing rate, the strength required to support the closure, the risk of infection, and cosmetic considerations.

\section*{1. Classification by Absorption Profile}

\subsection*{Absorbable Sutures}
These sutures are degraded by the body via enzymatic digestion or hydrolysis, losing their tensile strength over a predictable period.
\begin{itemize}
  \item \textbf{Surgical Gut (Plain or Chromic):}
  \begin{itemize}
    \item \textbf{Material:} Purified collagen from sheep/cattle intestine.
    \item \textbf{Profile:} Plain gut loses strength in 7-10 days (absorbed in 70 days); Chromic gut is treated with chromium salts to delay absorption to 14-21 days (absorbed in 90 days).
    \item \textbf{Common Uses:} Mucosal closures, superficial vascular ligation.
  \end{itemize}
  \item \textbf{Polyglactin 910 (Vicryl):}
  \begin{itemize}
    \item \textbf{Material:} Synthetic braided copolymer.
    \item \textbf{Profile:} Retains $50\%$ strength at 21 days; absorbed via hydrolysis in 56-70 days.
    \item \textbf{Common Uses:} Subcutaneous tissue approximation, bowel anastomosis.
  \end{itemize}
  \item \textbf{Poliglecaprone 25 (Monocryl):}
  \begin{itemize}
    \item \textbf{Material:} Synthetic monofilament copolymer.
    \item \textbf{Profile:} Retains $50\%$ strength at 7-14 days; absorbed in 90-120 days.
    \item \textbf{Common Uses:} Subcuticular skin closure.
  \end{itemize}
  \item \textbf{Polydioxanone (PDS II):}
  \begin{itemize}
    \item \textbf{Material:} Synthetic monofilament polymer.
    \item \textbf{Profile:} Retains $50\%$ strength at 28-42 days; absorbed in 180-210 days.
    \item \textbf{Common Uses:} Fascial closures, pediatric cardiovascular surgeries.
  \end{itemize}
\end{itemize}

\subsection*{Non-Absorbable Sutures}
These sutures resist degradation and remain permanently in the body (or are removed postoperatively if used on the skin).
\begin{itemize}
  \item \textbf{Surgical Silk:}
  \begin{itemize}
    \item \textbf{Material:} Natural protein filament.
    \item \textbf{Profile:} High tissue reactivity; degrades slowly over years.
    \item \textbf{Common Uses:} Secure drain fixation, temporary surgical retraction.
  \end{itemize}
  \item \textbf{Nylon (Ethilon, Dermalon):}
  \begin{itemize}
    \item \textbf{Material:} Synthetic monofilament polyamide.
    \item \textbf{Profile:} Minimal tissue reaction; high memory (requires multiple knots).
    \item \textbf{Common Uses:} Cutaneous skin closure.
  \end{itemize}
  \item \textbf{Polypropylene (Prolene):}
  \begin{itemize}
    \item \textbf{Material:} Synthetic monofilament.
    \item \textbf{Profile:} Extremely low tissue reactivity; high plasticity; retains strength indefinitely.
    \item \textbf{Common Uses:} Vascular anastomoses, fascial closures (hernia repair).
  \end{itemize}
  \item \textbf{Polyester (Ethibond):}
  \begin{itemize}
    \item \textbf{Material:} Synthetic braided.
    \item \textbf{Profile:} Permanent strength; low reactivity.
    \item \textbf{Common Uses:} Valve replacements, tendon/ligament repairs.
  \end{itemize}
\end{itemize}

\section*{2. Monofilament vs. Braided (Multifilament)}
\begin{itemize}
  \item \textbf{Monofilament (e.g., Prolene, Nylon, Monocryl, PDS):}
  \begin{itemize}
    \item \textbf{Advantages:} Smooth passage through tissues, lower risk of harboring bacteria (no capillarity).
    \item \textbf{Disadvantages:} Subject to damage if clamped; high memory makes handling and knot tying more difficult.
  \end{itemize}
  \item \textbf{Braided (e.g., Silk, Vicryl, Ethibond):}
  \begin{itemize}
    \item \textbf{Advantages:} Excellent handling, low memory, high knot security.
    \item \textbf{Disadvantages:} Higher risk of tissue drag; can act as a wick to harbor bacteria (capillarity).
  \end{itemize}
\end{itemize}

\section*{3. Suture Sizing (USP Scale)}
Suture diameter is defined by the United States Pharmacopeia (USP) scale, where "0" is mid-size, whole numbers indicate increasing thickness, and numbers followed by "-0" indicate decreasing thickness.
\begin{itemize}
  \item \textbf{Size 1 and 2-0:} Heavy fascial closures, bone anchor fixation.
  \item \textbf{Size 3-0 and 4-0:} General subcutaneous tissues, bowel anastomoses, superficial skin closures.
  \item \textbf{Size 5-0 and 6-0:} Face/neck skin closures (reduces scarring), small blood vessel anastomoses.
  \item \textbf{Size 7-0 to 11-0:} Microsurgery, ophthalmic procedures, coronary artery anastomoses.
\end{itemize}

\clearpage
